\chapter{Initial configuration}
\label{chap:initial_configuration}

\section{Greenbone}
\label{sec:greenbone}
Greenbone/OpenVAS is an open-source vulnerability 
scanning tool that provides comprehensive security 
assessments for IT systems.
To set up Greenbone/OpenVAS, firstly download the official 
Greenbone/OpenVAS OVA file from the Greenbone website.\\
\href{https://www.greenbone.net/en/openvas-free/}{https://www.greenbone.net/en/openvas-free/} \\
Then, import the OVA file into
VirtualBox 6.1 or higher.
Assign at least 5GB of RAM and 2 CPU cores to the virtual machine. 
It's better to use about 12GB of RAM.
For the first configuration, use the bridge network adapter, which allows the virtual machine to be on the same network as the host machine.
Start the Greenbone/OpenVAS virtual machine and follow the on-screen instructions to complete the initial setup.
Assign a user for the web interface.
Once the setup is complete, access the Greenbone/OpenVAS update the feed 
by going to \texttt{Maintenance>Feed>Update} and wait it to synchronize
(ususally hours).

\section{Debian}
\label{sec:debian}
Debian is an open-source operating system based on GNU/Linux.
It will be used as a base to create a vulnerable machine.
\subsection{Apache}
\label{sec:apache}
Install Apache web server and create weak certificates for it.
configure Apache to allow the use of weak protocols 
and ciphers, such as SSLv3, TLS1.1, DES, RC4 and MD5.
\inputminted{bash}{../project/apache/setup.sh}
Any website can be used as exposed service. 
In our case we used an old project of ours,
properly modified to be vulnerable.\\
\texttt{\href{https://github.com/ozneroL541/artigianato-online.git}
{git clone -b vulnerable https://github.com/ozneroL541/artigianato-online.git}}
To run it execute the following commands in the terminal:
\begin{minted}{bash}
cd artigianato-online/src
docker compose up --remove-orphans -d
\end{minted}

\subsection{Cups}
\label{sec:cups}
Install the CUPS printing system and configure it to allow remote administration without authentication.
\inputminted{bash}{../project/cups/setup.sh}

%\subsection{Docker Vulnet}
%\label{sec:docker-vulnet}

